\begin{longtable}{@{}p{4.4cm}p{4.2cm}p{5.4cm}@{}}
\caption{Pemetaan variabel dataset ke topik pedoman.}
\label{tab:pemetaan-variabel} \\
\toprule
\textbf{Variabel pada dataset} & \textbf{Kondisi yang dapat dipicu} & \textbf{Topik pedoman yang diperlukan} \\
\midrule
\endfirsthead

\multicolumn{3}{@{}l}{\small\emph{Tabel \thetable{} (lanjutan)}} \\
\toprule
\textbf{Variabel pada dataset} & \textbf{Kondisi yang dapat dipicu} & \textbf{Topik pedoman yang diperlukan} \\
\midrule
\endhead

\bottomrule
\endlastfoot

\texttt{glucose\_level}, \texttt{finger\_stick} & Kadar glukosa dalam atau di luar rentang sasaran & Pemantauan dan target glikemik \\[2pt]

\texttt{finger\_stick} rendah, \texttt{hypo\_event} & Hipoglikemia & Tata laksana hipoglikemia \\[2pt]

\texttt{glucose\_level} tinggi berkepanjangan & Hiperglikemia & Tata laksana hiperglikemia \\[2pt]

\texttt{glucose\_level} sangat tinggi & Risiko ketoasidosis & Ketoasidosis diabetik \\[2pt]

\texttt{bolus}, \texttt{basal}, \texttt{temp\_basal} & Penyesuaian dosis insulin & Terapi insulin \\[2pt]

\texttt{meal} beserta estimasi karbohidrat & Asupan karbohidrat & Nutrisi dan penghitungan karbohidrat \\[2pt]

\texttt{exercise}, \texttt{work} beserta intensitas & Aktivitas fisik & Pengelolaan saat aktivitas fisik \\[2pt]

\texttt{illness} & Kondisi sakit penyerta & Pengelolaan pada hari sakit \\[2pt]

\texttt{stressors} & Stres psikologis & Dipertahankan sebagai variabel \emph{logbook}, tidak dijadikan pemicu penelusuran karena jumlah kejadiannya terlalu sedikit \\

\end{longtable}

\noindent\footnotesize Pemetaan disusun melalui penelaahan kepadatan istilah kunci pada setiap dokumen korpus, kemudian divalidasi secara manual. Atribut durasi pada kanal \texttt{exercise} tidak diekstraksi, sehingga aktivitas fisik terwakili hanya oleh skor intensitas.\normalsize
